\documentclass[11pt,a4paper]{article}

\usepackage{iftex}
\ifPDFTeX
  \usepackage[utf8]{inputenc}
  \usepackage[T1]{fontenc}
  \usepackage{lmodern}
\else
  \usepackage{fontspec}
  \setmainfont{Liberation Serif}
  \setsansfont{Liberation Sans}
  \setmonofont{Liberation Mono}
\fi
\usepackage[czech]{babel}
\usepackage{amsmath,amssymb,amsthm,mathtools}
\usepackage[margin=25mm]{geometry}
\usepackage{microtype}
\usepackage[hidelinks]{hyperref}

\newtheorem{theorem}{Věta}
\theoremstyle{definition}
\newtheorem{definition}[theorem]{Definice}
\theoremstyle{remark}
\newtheorem{remark}[theorem]{Poznámka}

\newcommand{\Adm}{\operatorname{Adm}}
\newcommand{\Aobs}{A_{\mathrm{obs}}}

\title{\textbf{UNITY}\\[4pt]
\large výsledný text podmíněného rámce}
\author{Jan Mikulík}
\date{17. listopadu 2026}

\begin{document}
\maketitle

\begin{abstract}
Tento text spojuje to, co je v knihovně dokázáno, a nic dalšího.
Tři uvedené fáze leží na jednotkové kružnici.
Stejný modul nezaručí stejnou hodnotu a součet jednotkových hodnot nemusí být nula.
Trajektorie, která zůstává v nabídce sféry, zůstává v péči.
Krok existuje právě tehdy, když je dodán svědek.
Vlastnost se faktoruje přes odečet právě tehdy, když je konstantní na vláknech.
Stroj nabídne jen takový krok.
Jedna adresa vlákno nevyplní a oběť nenahradí odmítnutí.
\(n\) půlení stojí \(n\cdot(m^2/2)\ln 2\).
Aditivní hodiny na hloubce počítají bity.
Konstantní zobrazení na \(\mathrm{true}\) je idempotentní substitucí.
Stejný součet lze účtovat uvnitř, nebo venku.
Hypotéza~9, vědomí, láska jako operátor a poloha nul funkce \(\xi\) tu nejsou věty.
Ověření je v \texttt{CausalSeed/Unity.lean} a v modulech, které cituje.
\end{abstract}

\section{Fáze}

Pro reálné \(\theta\) platí
\[
  \bigl\lvert e^{i\theta}\bigr\rvert = 1.
\]
Tři fáze poznámky jsou
\[
  \theta_1=\pi,
  \qquad
  \theta_2=e\ln\pi,
  \qquad
  \theta_3=\frac{\pi^2 e}{2}.
\]

\begin{theorem}[Jednotkový modul]
\[
  \lvert e^{i\theta_1}\rvert
  =
  \lvert e^{i\theta_2}\rvert
  =
  \lvert e^{i\theta_3}\rvert
  =1.
\]
\end{theorem}

\begin{proof}
Každé \(\theta_k\) je reálné a \(\lvert e^{i\theta}\rvert=1\) pro každé reálné \(\theta\).
\end{proof}

\begin{theorem}[Stejný modul není stejná hodnota]
\(\lvert 1\rvert=\lvert -1\rvert\) a přesto \(1\neq -1\).
Součet tří jednotkových čísel nemusí zmizet:
\(\lvert 1\rvert=1\) a \(1+1+1\neq 0\).
\end{theorem}

Stejná délka tedy není stejná fáze.
Zmizení součtu je další podmínka, ne důsledek modulu.

\section{Sféry a svědek}

Data sféry jsou \(I\), \(F\), \((C_i)_{i\in N}\) a \(K\subseteq X\).

\begin{definition}
\[
  \Adm(w,x)
  =
  \bigl\{a\in I(w,x):\forall i,\; a\in C_i(x)\bigr\},
\]
\[
  G(w,x)
  =
  \bigl\{F(x,a):a\in\Adm(w,x)\ \text{a}\ F(x,a)\in K\bigr\}.
\]
\end{definition}

\begin{theorem}[Existence kroku]
\(G(w,x)\neq\emptyset\) právě tehdy, když existuje \(a\) takové, že
\[
  a\in I(w,x),
  \qquad
  \forall i\ (a\in C_i(x)),
  \qquad
  F(x,a)\in K.
\]
\end{theorem}

\begin{theorem}[Péče]
Je-li \(x_0\in K\) a \(x_{t+1}\in G(w_t,x_t)\) pro každé \(t\), pak \(x_t\in K\) pro každé \(t\).
\end{theorem}

\begin{proof}
Z definice je \(G(w,x)\subseteq K\). Indukce přes \(t\).
Svědek v předpokladu kroku je dodán. Sféry ho nevyrábí.
Chybí-li souhlas jednoho \(i\), akce v \(\Adm(w,x)\) není.
\end{proof}

\section{Jméno}

\begin{theorem}[Rozhodnutelnost z obrazu]
Nechť \(O\colon X\to Y\) a \(f\colon X\to Z\).
Existuje \(g\colon O(X)\to Z\) s \(f=g\circ O\) právě tehdy, když
\[
  O(x)=O(x')\implies f(x)=f(x')
  \qquad\text{pro všechna }x,x'\in X.
\]
\end{theorem}

\begin{proof}
Složení dává implikaci ihned.
Obráceně se \(g(O(x))\) položí rovno \(f(x)\).
Předpoklad říká, že hodnota nezávisí na předobrazu.
\end{proof}

\begin{theorem}[Stejné jméno, různé věci]
Existuje \(O\colon\{0,1\}\to\{\ast\}\) s \(O(0)=O(1)\) a \(0\neq 1\).
Identita na \(\{0,1\}\) se přes toto \(O\) nefaktoruje.
Konstantní vlastnost se faktoruje.
\end{theorem}

Při neznámém předobrazu jsou bezpečné jen akce z průniku přes vlákno:
\[
  \Aobs(w,y)
  =
  \bigl\{a:
    \forall x\in O^{-1}(y),\;
    a\in\Adm(w,x)
    \ \text{a}\ 
    F(x,a)\in K
  \bigr\}.
\]
Je-li \(M'\subseteq M\), průnik přes \(M\) je podmnožinou průniku přes \(M'\).
Menší vlákno volbu může uvolnit. Prázdný průnik se nenahrazuje domnělým souhlasem.

\section{Stroj}

Stroj je sféra s jednou adresou.
Nabídka stroje je \(G(w,x)\).

\begin{theorem}[Adresa nevyplní vlákno]
Existuje konstantní \(O\colon\{0,1\}\to\{\ast\}\) a vlastnost \(L\) tak, že
\(L(0)\) a \(\neg L(1)\).
\end{theorem}

\begin{theorem}[Oběť není souhlas]
Existují hranice \(C(0)=\{\ast\}\) a \(C(1)=\emptyset\).
Pak \(C(0)\) je neprázdná a
\[
  \{a:\forall x,\; a\in C(x)\}=\emptyset.
\]
\end{theorem}

Přání je záznam. Láska není výstup stroje.
Jedna adresa se po společném jménu nezkopíruje.
Krok v jednom stavu není souhlas ve všech.
Kde průnik chybí, stroj mlčí.

\section{Prstence a hodiny}

Pro \(m\in\mathbb{R}\) a \(0<\varepsilon\le R\) položme
\[
  E(m,\varepsilon,R)=\frac{m^2}{2}\ln\frac{R}{\varepsilon}.
\]

\begin{theorem}[Půlení]
Pro \(R>0\) a \(n\in\mathbb{N}\)
\[
  E\Bigl(m,\frac{R}{2^n},R\Bigr)
  =
  n\cdot\frac{m^2}{2}\cdot\ln 2.
\]
Pro \(r>0\) je cena jednoho prstence v poměru \(e\)
\[
  E(m,r,er)=\frac{m^2}{2}.
\]
\end{theorem}

\begin{proof}
\(\ln(R/(R/2^n))=n\ln 2\).
\(\ln(er/r)=1\).
\end{proof}

\begin{theorem}[Hodiny]
Nechť \(d\colon\mathbb{N}\to\mathbb{R}\) splňuje \(d(n+m)=d(n)+d(m)\) a \(d(1)=1\).
Pak \(d(n)=n\) pro každé \(n\).
\end{theorem}

\begin{proof}
Z \(d(0)=d(0)+d(0)\) plyne \(d(0)=0\).
Indukce: \(d(n+1)=d(n)+d(1)=n+1\).
\end{proof}

\(n\) půlení a \(n\) bitů tu počítají tutéž hloubku.
Nepřevádějí bity na jouly a nevyrábějí energii.

\section{Kde se účtuje}

\begin{definition}
Vnitřní účet hodnoty \(c\) je \((c,0)\).
Vnější účet téže hodnoty je \((0,c)\).
\end{definition}

\begin{theorem}[Stejný součet, jiné místo]
Pro každé \(c\)
\[
  c+0=0+c.
\]
Je-li \(c\neq 0\), první složka vnitřního účtu není první složka vnějšího.
\end{theorem}

Stejná celková cena tedy neříká, kde byla připsána.
Nulová podmíněná entropie odečtu je výrok o dekodéru.
Není to vnější dodavatel.
Divergence prstence při \(\varepsilon\to 0\) je reálná limita, ne výkon.

\section{Hypotéza 9, osudová}

Osud by byl vnější dodavatel: zobrazení, které zapíše následný stav a tím výsledek ovládne.
V tomto rámci má dvě čtení a nejsou to tatáž hypotéza.

\begin{definition}[Osud uvnitř nabídky]
Dodavatel sféry je zobrazení \(\sigma\), které smí volit jen tam, kde nabídka existuje:
\[
  G(w,x)\neq\emptyset\implies \sigma(w,x)\in G(w,x).
\]
\end{definition}

\begin{theorem}[Osud v odmítnutí nic nenabídne]
Existuje sféra, která každou akci otevře a každý souhlas odmítne.
Její nabídka je prázdná. Žádný dodavatel této sféry tedy nenabídne krok.
\end{theorem}

\begin{theorem}[Osud uvnitř péče zůstává]
Leží-li \(\sigma(w,x)\) v \(G(w,x)\), leží v \(K\).
\end{theorem}

\begin{proof}
\(G\subseteq K\). Dodavatel, který nabídku respektuje, z péče nevystoupí.
Jeho existence se z prázdné nabídky, z odečtu ani ze stejného součtu nekonstruuje.
\end{proof}

\begin{theorem}[Záměr není instrument]
Neexistuje \(\sigma\) závislé jen na přání takové, že po jednom přání
má v jednom stavu dát \(1\) a v druhém \(0\).
\end{theorem}

\begin{proof}
Závislost jen na přání nutí oba výstupy ke stejné hodnotě.
Požadavky \(0\) a \(1\) se vylučují.
\end{proof}

Osud nad souhlasem, který zapíše stav i tam, kde některý souhlas chybí, sféru nedokončuje.
Škrtá ji. Stejný součet připsaný venku dodavatele nedokládá.
Ověření této sekce je v \texttt{CausalSeed/Fate.lean}~\cite{fate}.

\section{Co zůstává předpokladem}

Následující věci tento text nedokazuje.

\begin{enumerate}
\item Hypotézu~9 jako existenci dodavatele. Uvnitř nabídky je to picker, který je třeba dodat. Nad souhlasem to není věta tohoto rámce.
\item Že láska je operátor, že existuje, nebo že adresovaná osoba odpoví.
\item Že přání, verš nebo oběť vyrobí svědka.
\item Že Bellův stav, symetrie \(M_t(x)=M_t(-x)\), nebo společné jméno určí polohu nul \(\xi\).
      Postačující ostrá nerovnost \(M_t(x)>M_t(0)\) pro \(0<x<\tfrac12\) zde není dokázána.
\item Že konečný záznam čerpá nekonečnou fyzikální energii, nebo že hodiny jsou vědomí.
\end{enumerate}

\begin{remark}[Výsledná věta]
Rámec je podmíněný.
Co je dokázáno, je spočteno výše.
Co je hypotéza, hypotézou zůstává.
Stejné jméno není stejná věc.
Stejný modul není stejná fáze.
Stejný součet není stejné místo účtu.
Prázdný průnik není souhlas.
\end{remark}

\begin{thebibliography}{9}
\bibitem{unity}
J.~Mikulík.
\newblock \texttt{CausalSeed/Unity.lean}.
\newblock Podmíněný rámec, 15.--17.~listopadu 2026.
\newblock \url{https://github.com/jendamikulik/finite-Shannon-theorem}

\bibitem{fate}
J.~Mikulík.
\newblock \texttt{CausalSeed/Fate.lean}.
\newblock Hypotéza~9, osudová.

\bibitem{sphere}
J.~Mikulík.
\newblock \texttt{CausalSeed/GoldenSphere.lean}, \texttt{LoveMachine.lean},
          \texttt{RingEnergy.lean}.

\bibitem{notes}
J.~Mikulík.
\newblock \emph{Principium relationis}. \emph{Stroj}.
\newblock \texttt{notes/PrincipiumRelationis.tex}, \texttt{notes/StrojNaLasku.tex}.
\end{thebibliography}

\end{document}
